% ---------------------------------------------------------------------
\section{Deney Düzeneği}\label{sec:setup}
% ---------------------------------------------------------------------

\subsection{Simülatör ve senaryolar}\label{sec:scenarios}

Tüm deneylerde, kaynak kodunda değişiklik yapılmadan \texttt{55ef368c} commit'indeki EV2Gym \cite{orfanoudakis2025ev2gym} kullanılmaktadır.
Bir sarmalayıcı (wrapper) gözlemleri hesaplar; ortak kabul edilen gereksinimleri, güvenlik katmanını, yedek kuralı ve eyleyici davranışını uygular ve gelecekteki varışları okumadan tüm ölçütleri kaydeder.
Her istasyonda tek bir trafonun arkasında $N$ şarj ünitesi bulunur; EV'ler simülatörün V2G destekli sekiz EV modelinden, her model için sabit bir verimle (bir model için 0,93, diğerleri için 0,90) çekilir.
İstenen enerji bataryanın tam dolu hâlidir; fiyatlar EV2Gym ile birlikte sunulan Hollanda gün öncesi fiyatlarıdır ve deşarj fiyatı şarj fiyatına eşittir.
Tablo~\ref{S-tab:scenarios} tüm senaryo parametrelerini listelemektedir.

Şebeke büyüklükleri ve varış hızı $N$ ile doğrusal olarak ölçeklenir.
Özellikle trafo anma gücü $P_{\sm{tr}}=\rho_{\sm{tr}}N$'dir; burada EV2Gym'in 25 şarj üniteli V2G yapılandırmasındaki gibi şarj ünitesi başına $\rho_{\sm{tr}}=4$ kW'lık sıkı bir düzey ve 7 kW'lık orta bir düzey kullanılır.
Esnek olmayan yük ve PV, ortalaması $0.4P_{\sm{tr}}$ olan çarpanlarla ölçeklenir ve gözlemdeki tahminler standart sapması \%10 olan Gauss gürültüsü içerir.
EV2Gym'in işyeri, kamusal ve özel varış profilleri, gece boyu süren oturumların 15 dakikalık $T=112$ adımdan oluşan 28 saatlik ufka sığmasını sağlayan saatlerde başlar.

Eğitim, her yığında (batch) $N\in\{10,20\}$ değerlerini, üç profili ve iki anma düzeyini karıştırır.
Test ızgarası $N\in\{5,10,20,50,100,200\}$ değerlerini, her iki anma düzeyini ve tüm profilleri; buna ek olarak $N\in\{20,200\}$ için $\rho_{\sm{tr}}=3$ kW'lık bir stres düzeyini ve EV2Gym'in özgün 25 şarj üniteli yapılandırmasını, toplamda 43 yapılandırmayı kapsar.
Dolayısıyla sıfır-atış ifadesi istasyon boyutunu ve stres düzeyini ifade eder; profiller ise eğitimde görülmüştür.
Her yapılandırma, tüm yöntemlerce paylaşılan ve $N\le50$ için 100, $N\ge100$ için 50 sabit test epizoduna sahiptir; böylece her politika 3.550 test epizodunda çalıştırılır.
Stres düzeyi, çalıştırılan tek duyarlılık çalışmasıdır.

\subsection{Karşılaştırılan yöntemler}\label{sec:baselines}

Her yöntem ve her değerlendirme geçişi aynı kabul edilen gereksinimleri alır ve eksik enerji her zaman bu gereksinimlere göre ölçülür.
Katman içermeyen yöntemler, $l_i$ ya da trafo bağlantısı olmaksızın yalnızca fiziksel kırpma $\clip(x_i,\pmn_i,u_i)$ işlemini uygular.

Öğrenen yöntemler arasında \method{} ve Deep Sets+safety layer, güvenlik katmanına sahip iki boyuttan bağımsız küme politikasıdır.
Attention+penalty, \method{} yönteminin cezaya dayalı ikizidir: güvenlik katmanı olmaksızın, kârdan eksik enerji ve aşırı yük enerjisinin $w$ ağırlığıyla çarpımı çıkarılarak elde edilen değerden öğrenir.
Bunun kısıtlı karşılığı olan Attention+PPO-Lagrangian, sabit ağırlığın yerine maliyet limiti sıfır olan Lagrange çarpanlarını kullanır \cite{altman2021constrained,tessler2018reward}.
$l_i=\pmn_i$ alındığında ``Attention+safety layer without departure bound'' (ayrılış sınırı olmayan dikkat+güvenlik katmanı) varyantı, katmanı \cite{wan2025logicinformed} çalışmasındaki gibi adım başına bir trafo izdüşümüne indirger ve ayrılış gereksinimini ceza içeren ödülden öğrenir.
Ceza ile eğitilmiş kontrol noktaları ayrıca güvenlik katmanı yalnızca kullanım aşamasında uygulanarak çalıştırılır; \method{} ve Attention+penalty ayrıca kabul kontrolü olmaksızın, varışta sınırlandırılmış talep korunarak çalıştırılır.
20 yuvaya kadar sıfırla doldurulmuş belirteçler üzerinde çalışan, katmanlı ve katmansız çok katmanlı algılayıcı (MLP), yalnızca $N\le20$ için değerlendirilen sabit boyutlu bir referanstır.
Dağılım içi bir referans olarak \method{} ayrıca $N=50$ için eğitilir ve $N\in\{50,100,200\}$ için değerlendirilir.

Karşılaştırmayı, her biri katmanlı ve katmansız olmak üzere altı kural tamamlar: olabildiğince hızlı (AFAP), olabildiğince geç (ALAP), döngüsel sıralama (round robin), LLF, sLLF \cite{chen2022smoothed} ve bir fiyat eşiği kuralı.
EV2Gym'in ekonomik MPC'sinin \cite{diazlondono2026open} HiGHS ile yaptığımız yeniden gerçeklemesi eMPC-HiGHS, her $N$ için 10 adımlık bir ufukla çalışır ve bu ufuk içindeki kesin fiyatları ve varışları bilir.
Konformal varyant \method{}-CP, $\alpha=0.1$ için bir kez eğitilir ve her $\alpha\in\{0.01,0.05,0.1,0.2\}$ için yeniden kalibre edilmiş sınırlarla değerlendirilir; ek bir geçiş olarak kalibre edilmemiş bir nokta tahmini de kullanılır.
EV-GNN \cite{orfanoudakis2025scalable} harici bir karşılaştırma yöntemi olarak planlanmış ancak çalıştırılmamıştır; GNN-DT \cite{orfanoudakis2026graph} ise her senaryo için bir optimizasyon çözücüsünün ürettiği veri kümeleri üzerinde çevrim dışı eğitildiğinden uyarlanmamıştır.
Planlanan ablasyonlar, ceza ağırlığı duyarlılıkları, diğer MPC varyantları ve diğer duyarlılık çalışmaları da çalıştırılmamıştır.

\subsection{Ölçütler}\label{sec:metrics}

Ana ölçüt, her yöntem için geçerli olan \Eq{eq:gap} ile tanımlı mutlak açıktır.
Cezalı kârlar, her yöntem ve geçiş için ortak kabul edilen gereksinime göre karşılanmayan enerjiyi kullanır; kâr, karşılanmayan enerji ve trafo aşırı yük enerjisi ayrıca raporlanır.
Güvenlik katmanı için kalış süresi boyunca olursuz adımla karşılaşmayan EV'lerin payı (E4), yedek kural durumlarına göre olursuz adımlar ve Lemma~\ref{S-lem:deficit} ile verilen açık sınırıyla karşılaştırmalı olarak ölçülen eksik enerji raporlanır.
Kayıtlar Teorem~\ref{thm:invariance}(b) hipotezi için EV başına bir bayrak taşımadığından, teoremin ayrılışa ilişkin kısmı epizot düzeyinde denetlenmektedir.

Garantinin maliyeti, aynı epizotlar üzerinde cezaya dayalı ikiz ile güvenlik katmanlı politika arasındaki farktır; birincil ölçü olarak $J^{\sm{pen}}$, ikincil ölçü olarak kâr kullanılır.
Bu maliyet, eşleştirilmiş epizotlar üzerinden toplamların oranıdır ve kabul edilen EV başına \texteuro{} ile $|J^{\sm{LP,ceil}}|$'nin yüzdesi olarak ifade edilir.
Pozitif bir değer, cezaya dayalı ikizin daha yüksek puan aldığı anlamına gelir.
\method{} yönteminin $N=50$ için eğitilen referansa göre aktarım açığı, aktarım kaybını yöntem kaybından ayırır.

\subsection{Ayarlama, istatistik ve hesaplama}\label{sec:hyper}

Ayarlanan yedi öğrenen yöntemin her biri, öğrenme oranı, mini yığın boyutu ve entropi katsayısı üzerinde tanımlı bir ızgaradan 12 yapılandırma almıştır.
Ceza içeren ödüller ızgaraya $w\in\{1,10,100\}$ ağırlığını, Lagrange yöntemi ise çarpan adım büyüklüğünü ekler.
Her yapılandırma eğitim karışımı üzerinde $0.5\times10^{6}$ adım eğitilmiş ve test tekrarlarından ayrık 24 doğrulama epizodunda değerlendirilmiştir; seçimde her yöntem için $w=10$ ile hesaplanan doğrulama $J^{\sm{pen}}$ değeri kullanılmıştır.
Seçilen yapılandırma iki tohumla $3\times10^{6}$ adım boyunca yeniden eğitilmiştir; Tablo~\ref{S-tab:hyper} hiperparametreleri listelemektedir.

İki tohumla tohumlar arası değişkenlik test edilemediğinden analiz, tohum ortalamalı politikaları ortak test tekrarları üzerinde karşılaştırmaktadır.
Tablolar, tohum ortalamalı politikanın ortalamasını iki tohum üzerindeki aralıkla birlikte raporlamaktadır.
Tekrar düzeyindeki farklar üzerinde eşleştirilmiş Wilcoxon işaretli sıralar testleri \cite{wilcoxon1945individual}, her istasyon boyutu için \method{} yöntemini diğer her yöntemle karşılaştırır.
Holm düzeltmesi bir sonuç tablosunun tüm karşılaştırmaları üzerinden (aile başına en fazla 351) uygulanır ve sıra-çift serili korelasyonun (rank-biserial) işareti yönü verir.
Ortalamanın güven aralıkları, test tekrarlarını her yapılandırma içinde yeniden örnekleyen bir bootstrap ile elde edilir ve açığın çeyrekler arası ortalaması (IQM) eğitim boyutlarında raporlanır \cite{agarwal2021deep}.
Wilcoxon testi medyan kaymasıyla, güven aralığı ise ortalamayla ilgilidir; büyük cezalı birkaç epizot tek başına ortalamayı kaydırabilir.
İkisinin uyuşmadığı durumlarda her ikisi de raporlanmakta ve karşılaştırma sonuçsuz (inconclusive) olarak nitelendirilmektedir.
Az sayıda tohuma dayanan sonuçlar yanıltıcı olabilir \cite{henderson2018deep}; bu nedenle tohum sayısı bir sınırlılıktır (Bölüm~\ref{S-app:stats}).

Gerçekleme PyTorch \cite{paszke2019pytorch}, CleanRL tarzında \cite{huang2022cleanrl} tek dosyalık bir PPO ve Adam optimize edicisini \cite{kingma2015adam} kullanmakta; güvenlik katmanı çift duyarlıkla çalışmaktadır.
Tüm deneyler WSL2 Ubuntu 24.04 altında Python 3.13.15 ve PyTorch 2.14.1 ile bir AMD Ryzen 7 7700 işlemci ve bir NVIDIA GeForce RTX 5070 Ti (16 GB) üzerinde çalıştırılmıştır.
Hazırlık 1,6 saat, ayarlama 7,9 saat, eğitim 10,2 saat ve değerlendirme 2,5 saat sürmüş; toplamda 153.100 epizot çalıştırılmıştır (Bölüm~\ref{S-app:compute}).
Kod, yapılandırmalar ve analiz betikleri makaleyle birlikte yayımlanacaktır \TODO{depo bağlantısı}.
